\section{问题四：实时波动电价下的购电策略}

\subsection{问题分析}

问题四要求在外网电价实时波动的条件下，重新制定问题二与问题三中的购电与储能调度策略。与前几问不同，当天真实电价在制定购电计划时尚未完全获得，而实际费用又需按照实时电价结算，因此本问首先需要解决未来电价未知条件下的价格预测问题，并在此基础上重新进行购电优化。

针对图~\ref{fig:att4} 所示的日内峰谷结构与逐日价格波动，本文根据决策日前已经获得的历史电价，采用历史同刻均值、星期修正和近期修正相结合的方法预测当天各时段电价，以保证预测过程不使用未来信息。随后分别沿用问题二和问题三的模型，根据预测电价以及负载、光伏预测结果制定购电与储能调度计划。最后，利用附件~4 的实际电价计算全年实际运行费用，通过总费用、紧急购电量及不同策略间的费用差异评价购电策略。

\subsection{日前电价预测模型建立}

记第 $d$ 日时段 $t$ 的真实电价与预测电价为 $p_{d,t}$、$\hat p_{d,t}$。利用决策日前的历史数据，将预测分为同刻基础价格、星期修正与近期修正三部分：
\begin{equation}
  \hat p_{d,t}=b_{d,t}+\Delta_{\mathrm{week},d,t}+\Delta_{\mathrm{recent},d,t},
  \qquad t=1,\ldots,144.
  \label{eq:q4_price_forecast}
\end{equation}
设 $n_d$ 为此前历史日数，$w(d)$ 为星期编号，$\mathcal W_d=\{j<d\mid w(j)=w(d)\}$，$m_d=\min(7,n_d)$，则
\begin{align}
  b_{d,t}&=\frac{1}{n_d}\sum_{j<d}p_{j,t}, \label{eq:q4_price_base}\\
  \Delta_{\mathrm{week},d,t}&=\frac{1}{|\mathcal W_d|}\sum_{j\in\mathcal W_d}r_{j,t}, \label{eq:q4_week_correction}\\
  \Delta_{\mathrm{recent},d,t}&=\frac{1}{m_d}\sum_{k=1}^{m_d}e_{d-k,t}. \label{eq:q4_recent_correction}
\end{align}
其中，$r_{j,t}=p_{j,t}-b_{j,t}$、$e_{j,t}=p_{j,t}-b_{j,t}-\Delta_{\mathrm{week},j,t}$ 为各历史日当时形成的两层残差；无相应历史样本时，修正项取零。

$1$ 月 $1$ 日以附件~1 电价初始化 $b_{1,t}$，两项修正均取零；此后每日先预测，待真实电价获得后再更新历史价格与残差。$1$ 月仅用于积累历史，不计入正式费用；$2$ 月起继续按相同规则逐日更新。

\subsection{模型求解}

本问继承问题二、三的预测方法、调度参数及物理约束，将式~\eqref{eq:q2_lp}、\eqref{eq:q3_lp0}、\eqref{eq:q3_lp_adj} 中的决策电价替换为 $\hat p_{d,t}$。正式回放仍从 $2$ 月 $1$ 日 $6000$ kWh 初始储电量出发，$1$ 月储电量不滚入。

Q4-2 在每日 $0{:}00$ 制定合同并保持日内不变，按问题二的规则回放。Q4-3 在 $6/12/18$ 时点根据更新的负载、光伏信息调整合同，但电价始终采用当天 $0{:}00$ 的预测；次日展望段按相同日内时刻延伸该预测电价。

两模型均调用 CBC 逐日求解，各次优化均获得最优解。实际费用分别按式~\eqref{eq:q2_cost}、\eqref{eq:q3_cost} 结算，其中电价取附件~4 的真实值 $\pi_{d,t}=p_{d,t}$，从而区分决策所用预测价与实际结算价。

\subsection{求解结果}

通过上述预测与调度过程，分别得到实时波动电价下问题二的日前购电方案和问题三的日内滚动购电方案。下面分别对 Q4-2 和 Q4-3 的求解结果进行分析。

% 避免三级标题孤悬于页底。
\Needspace{3\baselineskip}
\subsubsection{Q4-2 日前购电结果}

Q4-2 日前购电方案全年总费用为 $14\,413\,479.02$ 元。为进一步分析该方案的全年运行特征，图~\ref{fig:q4_dispatch} 给出了日初储电量、日购电量、弃电量、日成本以及紧急购电情况。

\begin{figure}[!htbp]
  \centering
  \includegraphics[width=0.78\textwidth]{q4_official_dispatch_details.png}
  \caption{Q4-2 日前购电方案日级运行情况}
  \label{fig:q4_dispatch}
\end{figure}

由图~\ref{fig:q4_dispatch} 可见，Q4-2 日前购电方案的运行结果具有明显的季节变化。图(a)中日初实际储电量始终保持在允许范围内，并围绕计划日末储电量目标波动，表明储能能够保持跨日连续运行。图(b)中日购电量与弃电量随季节发生明显变化，反映了全年负载与光伏发电水平变化对外网购电需求的影响。

由图(c)可见，日成本主要由计划购电费用构成，紧急购电费用在大部分日期较低，仅在少数日期明显增加。图(d)进一步表明，发生紧急购电的日期中，应急购电时段数增加时，应急购电成本总体也随之上升。因此，Q4-2 日前购电方案能够满足全年大部分时段的供电需求，但由于全天计划在 $0{:}00$ 提前制定，当负载或光伏预测出现较大偏差时，少数日期仍需要通过紧急购电弥补供电缺口。


% 避免三级标题孤悬于页底。
\Needspace{3\baselineskip}
\subsubsection{Q4-3 日内滚动购电结果}

为分析日内新增信息对购电策略的改善作用，设置三种方案进行比较。其中，N0 为日前固定方案，即每日 $0{:}00$ 制定购电计划后不再进行日内调整；M0L 为选择性日内调整方案，仅在满足调整条件时修改后续尚未执行的购电计划；LA 为本文采用的日内滚动调整方案，在 $6{:}00$、$12{:}00$ 和 $18{:}00$ 根据新获得的负载和光伏信息重新调整后续购电计划。三种方案的全年费用与紧急购电情况见表~\ref{tab:q4_all}。

\begin{table}[!htbp]
  \centering
  \caption{波动电价下问题三的全年结果比较}
  \label{tab:q4_all}
  \begingroup
  \small
  \setlength{\tabcolsep}{5pt}
  \begin{tabular}{@{}lrrr@{}}
    \toprule
    方案 & 总费用（元） & 紧急购电量（kWh） & 较 N0 节省（元） \\
    \midrule
    N0（不调整）
      & $15\,983\,608.51$ & $466\,576.26$ & --- \\
    M0L（选择性调整）
      & $14\,826\,830.64$ & $218\,215.59$ & $1\,156\,777.87$ \\
    \textbf{LA（本文）}
      & \textbf{$14\,001\,199.26$} & \textbf{$67\,153.26$} & \textbf{$1\,982\,409.26$} \\
    \bottomrule
  \end{tabular}

  \par\smallskip
  \begin{minipage}{0.94\textwidth}
    \footnotesize\textit{注：}LA 的合同购电费用为 $12\,979\,047.00$ 元，
    紧急购电费用为 $370\,783.78$ 元，其余费用来自少买与多买偏差结算。
    各方案均采用附件~4 的真实电价按照相同规则计算实际费用。
  \end{minipage}
  \endgroup
\end{table}

由表~\ref{tab:q4_all} 可见，LA 的全年总费用为 $14\,001\,199.26$ 元，较不进行日内调整的 N0 减少 $1\,982\,409.26$ 元，降幅约为 $12.40\%$；紧急购电量由 $466\,576.26$ kWh 降至 $67\,153.26$ kWh，降幅约为 $85.6\%$。与 M0L 相比，LA 仍进一步节省 $825\,631.39$ 元，降幅约为 $5.57\%$。

结果表明，仅依靠 $0{:}00$ 制定的日前计划难以完全应对当天负载和光伏预测误差。LA 在 $6{:}00$、$12{:}00$ 和 $18{:}00$ 利用新获得的信息调整尚未执行的购电计划后，能够及时修正原计划与实际供需之间的偏差，从而显著减少高价紧急购电并降低全年总费用。

\begin{figure}[!htbp]
  \centering
  \begin{subfigure}[t]{0.96\textwidth}
    \vspace{0pt}\centering
    \includegraphics[width=\linewidth]{q4_cost_structure.pdf}
    \caption{全年成本结构}
    \label{fig:q4_cost}
  \end{subfigure}\\[-6pt]
  \begin{subfigure}[t]{0.96\textwidth}
    \vspace{0pt}\centering
    \includegraphics[width=\linewidth]{q4_daily_cost_distribution.pdf}
    \caption{日总成本分布}
    \label{fig:q4_daily}
  \end{subfigure}
  \caption{Q4-3 三种策略的费用结构与日费用分布}
  \label{fig:q4_cost_daily}
\end{figure}

由图~\ref{fig:q4_cost} 可见，三种策略的费用差异主要来自紧急购电。N0 的紧急购电费用较高，M0L 通过选择性调整有所降低，而 LA 进一步压低了紧急购电费用，因此全年总费用最低。图~\ref{fig:q4_daily} 进一步表明，LA 的日总费用整体较低，同时高费用日期明显减少。这说明日内滚动调整不仅能够降低全年总费用，还能够减少预测偏差造成的高费用情况。


% 避免三级标题孤悬于页底。
\Needspace{3\baselineskip}
\subsubsection{代表日期结果分析}

为进一步考察两类购电方案在不同季节下的运行结果，选取四个代表日期，其购电量、购电费用及紧急购电量见表~\ref{tab:q4_dates}。

\begin{table}[!htbp]
  \centering
  \caption{波动电价下指定日期的购电结果}
  \label{tab:q4_dates}
  \begingroup
  \small
  \setlength{\tabcolsep}{3pt}
  \begin{tabular}{@{}lrrr@{\hspace{7pt}}rrr@{}}
    \toprule
    日期 & \multicolumn{3}{c}{Q4-2 日前购电方案}
         & \multicolumn{3}{c}{Q4-3 日内滚动方案} \\
    \cmidrule(lr){2-4}\cmidrule(lr){5-7}
    & \shortstack{购电量\\（kWh）}
    & \shortstack{购电费\\（元）}
    & \shortstack{紧急购电量\\（kWh）}
    & \shortstack{购电量\\（kWh）}
    & \shortstack{购电费\\（元）}
    & \shortstack{紧急购电量\\（kWh）} \\
    \midrule
    2025-03-20 & $70\,570.86$ & $44\,640.38$ & 0.00
               & $66\,824.83$ & $44\,186.70$ & 165.68 \\
    2025-06-21 & $35\,205.93$ & $17\,564.67$ & 0.00
               & $33\,936.26$ & $16\,638.47$ & 0.00 \\
    2025-09-23 & $68\,848.07$ & $44\,716.63$ & 272.06
               & $67\,439.16$ & $47\,671.18$ & 631.22 \\
    2025-12-21 & $98\,970.30$ & $75\,354.93$ & 0.00
               & $96\,674.19$ & $75\,675.02$ & 291.73 \\
    \bottomrule
  \end{tabular}
  \endgroup
\end{table}

由表~\ref{tab:q4_dates} 可见，两类方案的购电需求均具有明显的季节差异。以 $6$ 月 $21$ 日和 $12$ 月 $21$ 日为例，Q4-2 的全天购电量分别为 $35\,205.93$ kWh 和 $98\,970.30$ kWh，Q4-3 分别为 $33\,936.26$ kWh 和 $96\,674.19$ kWh，冬季代表日的购电需求明显高于夏季代表日，反映了负载与光伏发电水平的季节变化。

同时，部分代表日期中 Q4-3 的紧急购电量高于 Q4-2。例如 $9$ 月 $23$ 日，两者分别为 $631.22$ kWh 和 $272.06$ kWh，说明日内滚动调整并不能保证每个单日的紧急购电量都低于日前方案。其效果应从全年范围评价：结合表~\ref{tab:q4_all} 可知，LA 相较不进行日内调整的 N0 将全年紧急购电量降低约 $85.6\%$，说明日内调整的主要优势在于整体降低预测偏差造成的紧急购电风险，而非保证所有单日结果均达到最低。

\subsection{模型检验}

为验证问题四调度结果的可靠性，从信息使用和运行约束两方面进行检验。对全年 $334$ 个决策日逐日核验后，未发现未来电价信息泄漏：每日 $0{:}00$ 的电价预测仅使用此前已获得的历史数据，$6{:}00$、$12{:}00$ 和 $18{:}00$ 的日内调整也未使用当天尚未获得的真实电价。

同时，将 Q4-2 与 Q4-3 的全年调度结果重新代入模型进行检查，所有决策日均满足储电量范围 $1200\le E_t\le10800$，且储能接口功率、电量平衡及不可向外网售电等约束均未出现违反情况。按照附件~4 的真实电价重新计算各项费用后，结果与模型输出一致。综上，所得方案符合题目规定的信息条件与运行约束，费用计算结果具有良好的数值一致性。
